\section{Background}
\label{sec:background}

\subsection{Duplex UMI sequencing}
\label{sec:bg-duplex}

Standard next-generation sequencing produces per-base error rates of $10^{-2}$ to $10^{-3}$, set by the accuracy of the sequencing chemistry.
Detecting a variant at 0.1\% VAF requires distinguishing a true signal from errors that occur at 10 to 100 times the signal frequency.
Duplex UMI sequencing solves this through two layers of error correction\autocite{Schmitt12duplex}.

Before PCR amplification, each original DNA molecule receives a pair of UMI tags, one ligated to each strand.
All PCR copies of that molecule carry the same UMI pair.
After sequencing, reads sharing a UMI pair are grouped into a \emph{molecule family}.
The first error correction step collapses all reads from one strand of a family into a \emph{single-strand consensus sequence} (SSC) by majority vote weighted by base quality.
Positions where most reads agree are called with high confidence; positions where reads disagree are flagged or masked.
This reduces the error rate to approximately $10^{-3}$ to $10^{-4}$, limited by systematic errors that affect all copies of one strand identically.

The second step compares the forward-strand SSC to the reverse-strand SSC.
A sequencing or PCR error on one strand is unlikely to produce the identical error on the complementary strand independently.
Positions where both SSCs agree form the \emph{duplex consensus sequence} (DCS), with an effective error rate of approximately $10^{-7}$: the product of the two independent per-strand error rates\autocite{Kennedy14duplex}.
This three-order-of-magnitude improvement over single-strand consensus is what makes sub-0.1\% VAF detection feasible.

The Twist Biosciences UMI adapter system\autocite{Twist22cfdna} uses the read structure \texttt{5M2S+T} on both R1 and R2: 5 bases of random UMI, 2 bases of skip sequence (a fixed monotemplate spacer from the ligation chemistry), and the remaining template sequence.
The 5-base UMI provides $4^5 = 1{,}024$ possible UMI values.
Duplex strand pairing uses a canonical form: given UMI pair $(A, B)$ from R1 and R2, the canonical identifier is $\min(A{+}B,\; B{+}A)$ by lexicographic order.
Reads with UMI pair $(A, B)$ and reads with UMI pair $(B, A)$ map to the same molecule on opposite strands.

The small UMI space creates a collision problem at high sequencing depth.
At 10{,}000$\times$ coverage, two unrelated molecules at the same genomic locus can share a UMI pair by chance.
Collisions merge distinct molecules into one family, masking real variants (false negatives) or averaging unrelated sequences (consensus errors).
Collision detection requires a secondary signal beyond the UMI itself, such as the genomic endpoints of the sequenced fragment.

\subsection{K-mer approaches to variant detection}
\label{sec:bg-kmer}

A k-mer is a subsequence of length $k$ extracted from a read or reference.
For $k = 31$, a 150-base read yields 120 overlapping 31-mers.
K-mer counting tools such as Jellyfish\autocite{Marcais11jellyfish} build a hash table mapping each distinct k-mer to its occurrence count across all input reads.
This index supports alignment-free analyses: any query about sequence content can be answered by looking up k-mer counts without aligning reads to a reference genome.

K-mer variant detection works by constructing a de Bruijn graph from the k-mer index and walking paths through it.
Each node in the graph is a k-mer; edges connect k-mers that overlap by $k{-}1$ bases.
Given a target reference sequence (for example, exon 7 of \emph{TP53}), the tool anchors the walk at the first and last k-mers of the reference, then enumerates all paths between these anchors that are supported by k-mer evidence in the index.
The reference path represents the wild-type allele.
Any alternative path represents a variant: a single-nucleotide substitution adds one divergent node, an insertion adds extra nodes, and a deletion skips nodes.
The ratio of variant-path k-mer counts to reference-path k-mer counts estimates the VAF.

km\autocite{Bourgey19km} implements this approach using Jellyfish k-mer databases.
It was designed for RNA-seq mutation detection and calls variants when the count ratio exceeds a user-specified threshold (the \texttt{-{}-ratio} parameter).
This threshold is a global cutoff with no statistical model: the same ratio at 100$\times$ and 10{,}000$\times$ coverage receives the same call, despite representing very different levels of evidence.
kmtools adds parallelisation (chunking targets across processes) and post-hoc filtering against known variant databases.

The fundamental limitation of this pipeline is that k-mer counts are integers with no provenance.
Jellyfish's canonical mode (\texttt{-C}) collapses forward and reverse-complement k-mers into a single count, erasing strand information entirely.
There is no per-k-mer record of how many molecules contributed the count, whether those molecules had duplex support, or what the base qualities were at the variant position.

\subsection{The information loss problem}
\label{sec:bg-information-loss}

The existing pipeline processes data through four stages, losing information at each transition.
Figure~\ref{fig:info-loss} illustrates the two pipelines side by side.

\begin{figure}[htbp]
\centering
\small
\begin{verbatim}
Existing pipeline:
  Raw reads --> HUMID (dedup) --> Jellyfish (count) --> km (call)
  [duplex       [one read        [integer              [ratio
   families,     per family,      counts,               threshold,
   strand info,  no consensus,    no strand,            no statistics,
   base quality] no metadata]     no molecules]         no confidence]

kam pipeline:
  Raw reads --> molecule assembly --> k-mer index --> variant calling
  [duplex       [SSC + DCS           [MoleculeEvidence  [molecule-count
   families,     consensus,           per k-mer:          statistics,
   strand info,  quality scores,      n_molecules,        duplex-aware
   base quality] family metadata]     n_duplex,           weighting,
                                      strand counts,      confidence
                                      error probs]        intervals]
\end{verbatim}
\caption{Information flow in the existing pipeline (top) and kam (bottom).
Each arrow represents a processing step.
Brackets list the information available at that stage.
The existing pipeline reduces duplex molecule families to integer k-mer counts by the second step.
kam preserves molecule identity and quality metadata through to variant calling.}
\label{fig:info-loss}
\end{figure}

Three specific losses matter for low-VAF variant detection.

\textbf{Loss of molecule identity.}
HUMID selects one representative read per UMI group and discards the rest.
Jellyfish then counts k-mers across all surviving reads.
A k-mer count of 500 could represent 500 independent molecules (strong evidence) or 5 molecules with 100 PCR copies each (weak evidence).
The distinction is critical at 0.1\% VAF, where a true variant may be supported by only 5 to 10 molecules out of 5{,}000 to 10{,}000 total.

\textbf{Loss of strand information.}
Jellyfish's canonical k-mer mode merges forward and reverse-complement k-mers.
A variant supported by molecules on only one strand is a hallmark of systematic artefact (oxidative damage, deamination), not a true somatic mutation.
Without strand information, the variant caller cannot apply strand bias filters.

\textbf{Loss of error correction quality.}
Duplex consensus bases have error rates of ${\sim}10^{-7}$.
Simplex consensus bases have error rates of ${\sim}10^{-3}$.
Singleton reads have error rates of ${\sim}10^{-2}$.
When all are collapsed into a single k-mer count, the variant caller treats evidence from a duplex-confirmed molecule identically to evidence from an uncorrected singleton.
A principled statistical model should weight these differently.

kam addresses all three losses by storing \texttt{MoleculeEvidence} per k-mer: distinct molecule count, duplex molecule count, per-strand simplex counts, and aggregate error probabilities.
The variant caller uses these fields directly, calling variants based on the number of independent molecules (not reads) and weighting duplex evidence above simplex evidence.
